\documentclass[11pt]{article}

\usepackage[margin=1in]{geometry}
\usepackage{graphicx}
\usepackage{amsmath}
\usepackage{booktabs}
\usepackage[round,authoryear]{natbib}
\usepackage[colorlinks=true,linkcolor=blue,citecolor=blue,urlcolor=blue]{hyperref}

\newcommand{\Rocrit}{\mathrm{Ro_{crit}}}
\newcommand{\Rosun}{\mathrm{Ro_\odot}}
\newcommand{\Msun}{M_\odot}

\title{Weakened Magnetic Braking in a Coupled Rotation \\ and XUV Evolution Model}
\author{Rory Barnes}
\date{28 August 2026}

\begin{document}
\maketitle

\begin{abstract}
\noindent
We have implemented the weakened magnetic braking (WMB) prescription of
Metcalfe et al.\ (2025) in a rotation-plus-XUV evolution model for
solar-type stars, applying the dynamo suppression factor simultaneously to
the wind braking torque and to the XUV luminosity. Using standard stellar
structure and torque prescriptions, the model without WMB reproduces the
Rossby number that Metcalfe et al.\ predict for a solar model at the
main-sequence turnoff, $\mathrm{Ro} \approx 1.8~\Rosun$. Adding WMB brings the
turnoff Rossby number down to their predicted $\approx 1.1~\Rosun$, but only if
the suppression factor is normalized to unity at small Rossby number. Applied
with the published normalization constant $N = 0.354$ taken literally, the
suppression reduces the torque by a factor of three from the zero-age main
sequence onward; the star then spins down too slowly ever to approach
$\Rocrit$, reaching only $\mathrm{Ro} = 0.85~\Rosun$ at 10~Gyr and arriving at
the solar age with a rotation period of 14.9~d rather than 25.4~d. We
recommend the normalized form and flag the choice for discussion.
\end{abstract}

\section{Purpose}

Metcalfe et al.\ (2025) show that the wind braking torque of solar-type stars
falls abruptly below standard predictions as the Rossby number approaches a
critical value for dynamo excitation, and that the same transition appears in
the large-scale magnetic field and in the X-ray luminosity. Their analysis is
of a snapshot sample: each star contributes one measured torque at one measured
Rossby number.

This note asks what happens when that prescription is embedded in a
time-dependent model and integrated for 10~Gyr. Two things motivate the
exercise. First, the suppression is self-referential --- it depends on the
Rossby number, which depends on the rotation period, which is set by the
suppressed torque --- so its consequences over an evolutionary track are not
obvious from the snapshot fit. Second, if the reduced X-ray luminosity reflects
reduced coronal heating rather than an artifact of the mass-loss calibration,
then WMB should also depress the XUV output that drives atmospheric escape from
any orbiting planets, which is the science this is ultimately aimed at.

\section{The model}

We evolve a $1~\Msun$ star from an age of 5~Myr to 10~Gyr, starting from a
rotation period of 1~d.

\subsection{Structure}

Stellar radius, luminosity, effective temperature, and radius of gyration come
from the \citet{Baraffe15} evolutionary grid, interpolated in mass and age. The
moment of inertia follows from the radius and radius of gyration, and its
change with time contributes to the rotational evolution through angular
momentum conservation, alongside the wind torque.

\subsection{Rossby number}

The convective turnover timescale is the \citet{CranmerSaar11} fit to effective
temperature,
%
\begin{equation}
  \tau_{\rm c}(T_{\rm eff}) = 314.24\,
    \exp\!\left[-\frac{T_{\rm eff}}{1952.5~\mathrm{K}}
      - \left(\frac{T_{\rm eff}}{6250~\mathrm{K}}\right)^{18}\right]
    + 0.002~\mathrm{d},
\end{equation}
%
and the Rossby number is $\mathrm{Ro} = P_{\rm rot}/\tau_{\rm c}$. On this
scale the solar value is $\Rosun = 2.168$, adopting the solar rotation period
used by \citet{Matt15}.

All Rossby numbers in this note are quoted in units of $\Rosun$. This matters
for comparison with Metcalfe et al., who infer $\tau_{\rm c}$ one pressure
scale height above the convective boundary from their own model grids. Absolute
Rossby numbers computed on different definitions of $\tau_{\rm c}$ are not
interchangeable; only the ratio to the solar value is portable, and we use it
throughout. The residual systematic is that the \emph{shape} of
$\tau_{\rm c}(T_{\rm eff})$ differs between the two prescriptions. That
difference cancels for the Sun but not at other masses, so the agreement
reported below should not be assumed to hold across the full F--K range without
a check.

\subsection{Wind braking torque}

The torque is the \citet{Matt15} prescription,
%
\begin{equation}
  \mathcal{T} =
  \begin{cases}
    \mathcal{T}_0\, \chi^2 \left(\dfrac{\Omega}{\Omega_\odot}\right),
      & \mathrm{Ro} \le \Rosun/\chi \quad \text{(saturated)},\\[2.2ex]
    \mathcal{T}_0
      \left(\dfrac{\tau_{\rm c}}{\tau_{{\rm c},\odot}}\right)^{2}
      \left(\dfrac{\Omega}{\Omega_\odot}\right)^{3},
      & \mathrm{Ro} > \Rosun/\chi \quad \text{(unsaturated)},
  \end{cases}
\end{equation}
%
with
$\mathcal{T}_0 = 6.3\times10^{30}\,(R/R_\odot)^{3.1}(M/\Msun)^{1/2}$~erg
and a saturation threshold $\chi = 10$.

\subsection{XUV luminosity}

The unsuppressed XUV luminosity follows the saturated power law of
\citet{Ribas05},
%
\begin{equation}
  L_{\rm XUV}^{\rm std}(t) =
  \begin{cases}
    f_{\rm sat}\, L_{\rm bol}, & t < t_{\rm sat},\\[1.2ex]
    f_{\rm sat}\, L_{\rm bol} \left(t/t_{\rm sat}\right)^{-\beta},
      & t \ge t_{\rm sat},
  \end{cases}
  \label{eq:ribas}
\end{equation}
%
with $f_{\rm sat} = 10^{-3}$, $t_{\rm sat} = 0.1$~Gyr, and $\beta = 1.23$. Note
that this is a broadband XUV luminosity, roughly three orders of magnitude
above the solar X-ray luminosity; see \S\ref{sec:caveats}.

\subsection{Weakened magnetic braking}
\label{sec:wmb}

\citet{Cameron17} proposed that the global solar dynamo behaves as a weakly
nonlinear system in the vicinity of a supercritical Hopf bifurcation. The
control parameter is the dynamo number $D \sim \mathrm{Ro}^{-2}$, and the
system sustains a periodic solution only while $D$ exceeds a critical value.
Near threshold the amplitude of a dynamo-driven observable scales as
$(D - D_{\rm crit})^{1/2}$, which in Rossby space is
$(\Rocrit - \mathrm{Ro})^{1/2}$. \citet{Metcalfe25} adopt this form to describe
the onset of WMB and fit it to their sample. We write the suppression factor as
%
\begin{equation}
  \Psi(x) =
  \begin{cases}
    N\,(\Rocrit - x)^{p}, & x < \Rocrit,\\[1.2ex]
    0, & x \ge \Rocrit,
  \end{cases}
  \qquad x \equiv \mathrm{Ro}/\Rosun,
  \label{eq:psi}
\end{equation}
%
with $p = 1/2$ and, from the Metcalfe et al.\ iterative fit,
$\Rocrit = 1.014 \pm 0.026$ and $N = 0.354 \pm 0.077$.

Because Metcalfe et al.\ attribute the torque deficit to a decline in Poynting
flux that starves coronal heating and wind acceleration alike, we apply the
same factor to both channels:
%
\begin{equation}
  \mathcal{T} \rightarrow \Psi\,\mathcal{T},
  \qquad
  L_{\rm XUV} \rightarrow
    \max\!\left[\Psi^{\,n}\, L_{\rm XUV}^{\rm std},\;
      f_{\rm floor}\, L_{\rm XUV}^{\rm Ribas}\right],
  \label{eq:apply}
\end{equation}
%
with $n = 1$. The floor, $f_{\rm floor} = 0.04$, represents the basal coronal
emission that survives the collapse of the global dynamo: four stars in the
Metcalfe et al.\ sample lie above $\Rocrit$ yet have X-ray detections between
$10^{26}$ and $10^{27}$~erg~s$^{-1}$, and the paper itself notes that the solar
chromospheric emission is dominated by an unstructured $\sim\!170$~G field that
is largely unaffected when the $\sim\!1$~G dipole disappears. No floor is
applied to the torque.

For comparison we also run the earlier prescription of \citet{vanSaders16}, in
which braking simply ceases above a threshold Rossby number,
$\Psi = 1$ for $x \le 0.959$ and $\Psi = 0$ above it. The threshold in those
units is slightly below solar, where the Metcalfe et al.\ value is slightly
above.

\section{Results}

Figure~\ref{fig:rotation} shows the rotational evolution and
Figure~\ref{fig:xuv} the XUV evolution, for the fiducial model and for the
three WMB variants.

\begin{figure}[htb]
  \centering
  \includegraphics[width=0.72\textwidth]{wmbRotation.pdf}
  \caption{Rotation period of a $1~\Msun$ star. The fiducial model, with no
    weakened magnetic braking, is the thick grey line; the WMB tracks coincide
    with it until suppression sets in and are drawn on top of it. The solar
    symbol marks $P_{\rm rot} = 25.4$~d at 4.57~Gyr. The literal $N = 0.354$
    track departs from the fiducial immediately and never recovers.}
  \label{fig:rotation}
\end{figure}

\begin{figure}[htb]
  \centering
  \includegraphics[width=0.72\textwidth]{wmbXUV.pdf}
  \caption{XUV luminosity for the same four models. The flat tails in the two
    tracks that reach $\Rocrit$ are the basal floor of
    Equation~\ref{eq:apply}. No solar point is plotted: the quantity here is a
    broadband XUV luminosity, which is not directly comparable to the X-ray
    luminosities of Metcalfe et al.\ (\S\ref{sec:caveats}).}
  \label{fig:xuv}
\end{figure}

The fiducial model reaches $P_{\rm rot} = 29.4$~d at the solar age, about 15\%
longer than the Sun's. Both the van Saders cutoff and the normalized Metcalfe
suppression bring the star into better agreement, at 28.8~d and 24.1~d
respectively. The literal $N = 0.354$ form gives 14.9~d, low by a factor of
1.7.

The XUV tracks separate later than the rotation tracks in the two threshold
models, because the suppression is negligible until the star nears $\Rocrit$.
Once it crosses, the XUV drops onto the basal floor within roughly a
gigayear --- an order-of-magnitude decline concentrated in the last third of
the main sequence, which is the feature most relevant to planetary atmospheric
escape. The literal form instead depresses the XUV by a factor of three
throughout, producing no late-time transition at all.

\section{Why the suppression is renormalized}
\label{sec:renorm}

Equation~\ref{eq:psi} with the published constants gives
$\Psi \rightarrow N\,\Rocrit^{\,p} = 0.357$ as $x \rightarrow 0$. Taken at face
value, the torque and the XUV of even the most rapidly rotating stars are
suppressed by nearly a factor of three --- in precisely the regime where
Metcalfe et al.\ state that standard spin-down models reproduce the
observations. This suggests that the vertical axis of their Figure~2 carries a
normalization that absorbs $N$, and that $\Psi$ is intended as a departure from
the standard model rather than a rescaling of it. We therefore also run a
variant with
%
\begin{equation}
  N_{\rm norm} = \Rocrit^{-p} = 0.9931,
\end{equation}
%
which sets $\Psi(x\!=\!0) = 1$ exactly and leaves $\Rocrit$ and $p$ untouched.

The evolutionary calculation discriminates sharply between the two choices,
because Metcalfe et al.\ supply a prediction that a snapshot fit cannot test:
a solar model with WMB should reach $\mathrm{Ro} \approx 1.1~\Rosun$ at the
main-sequence turnoff, against $\approx 1.8~\Rosun$ for a standard spin-down
model. Table~\ref{tab:results} collects the comparison and
Figure~\ref{fig:rossby} shows the tracks.

\begin{table}[htb]
  \centering
  \caption{Solar-analog diagnostics. Rossby numbers are in units of $\Rosun$.
    Metcalfe et al.\ predict $\mathrm{Ro} \approx 1.1$ at the turnoff with WMB
    and $\approx 1.8$ without.}
  \label{tab:results}
  \begin{tabular}{lccc}
    \toprule
    Model & $P_{\rm rot}$ at 4.57~Gyr & Ro at 10~Gyr & $\Psi$ at 4.57~Gyr \\
    \midrule
    Fiducial, no WMB           & 29.4~d & 1.80 & 1.00 \\
    van Saders cutoff          & 28.8~d & 1.16 & 0.00 \\
    Metcalfe, $N = 0.354$      & 14.9~d & 0.85 & 0.25 \\
    Metcalfe, normalized       & 24.1~d & 1.08 & 0.45 \\
    \midrule
    Sun / prediction           & 25.4~d & 1.1 (WMB), 1.8 (std) & --- \\
    \bottomrule
  \end{tabular}
\end{table}

\begin{figure}[htb]
  \centering
  \includegraphics[width=0.66\textwidth]{wmbRossby.pdf}
  \caption{Rossby number (top) and suppression factor (bottom). The dotted line
    marks $\Rocrit$ and the solar symbol marks $\mathrm{Ro} = \Rosun$ at
    4.57~Gyr. The fiducial model arrives at $1.80~\Rosun$ at 10~Gyr, matching
    the standard-model prediction of Metcalfe et al.; the van Saders and
    normalized Metcalfe models arrive at $1.16$ and $1.08$, matching their WMB
    prediction. The literal $N = 0.354$ model reaches only $0.85$ and is still
    spinning down.}
  \label{fig:rossby}
\end{figure}

Three results follow. First, the fiducial model independently recovers
$\mathrm{Ro} = 1.80$ at 10~Gyr, which is reassuring: two quite different
stellar structure and torque prescriptions agree on the standard-model
endpoint, so a discrepancy in the WMB runs is unlikely to be an artifact of our
choice of ingredients. Second, both threshold models land near the predicted
WMB value, at 1.16 and 1.08. Third, the literal normalization does not: it
reaches only $\mathrm{Ro} = 0.85$, and its suppression factor at that point is
still 0.145, so it is nowhere near the transition the model exists to describe.

The failure is structural rather than a matter of tuning. Suppressing the
torque threefold from the zero-age main sequence leaves the star rotating too
rapidly, hence at too small a Rossby number, hence further from $\Rocrit$ than
a standard model of the same age. The suppression prevents the star from ever
reaching the regime in which the suppression is supposed to operate. No choice
of $\Rocrit$ within the quoted uncertainty repairs this, because the problem
originates at small $x$, where the fit is least constrained by the data and
where the physical argument of \citet{Cameron17} says the dynamo should be
comfortably supercritical and therefore unsuppressed.

We therefore recommend the normalized form, and note that it is also the one
that puts the star closest to the solar rotation period at the solar age.

\section{Caveats}
\label{sec:caveats}

\paragraph{Bandpass.} Metcalfe et al.\ work in X-rays. The XUV luminosity here
is the broadband quantity of Equation~\ref{eq:ribas}, with
$L_{\rm XUV}/L_{\rm bol} = 10^{-3}$ at saturation against a present-day solar
$L_{\rm X}/L_{\rm bol} \approx 5\times10^{-7}$. The suppression factor is
dimensionless and multiplicative, so it transfers between bands without
rescaling, but the assumption that the EUV is suppressed by the same factor as
the X-rays is untested. If the EUV survives the dynamo collapse better than the
X-rays do, the consequences for planetary atmospheric escape are substantially
weaker than Figure~\ref{fig:xuv} implies.

\paragraph{Baseline.} The standard model against which Metcalfe et al.\ fit
$\Rocrit$ and $N$ is $L_{\rm X} \propto L_{\rm bol}/\mathrm{Ro}^{2}$. Our
unsuppressed XUV is instead the age power law of Equation~\ref{eq:ribas}, so
$\Psi$ multiplies a baseline different from the one it was calibrated against.
The two are not wildly different for a solar-mass star, since Ro and age are
themselves closely related along a track, but the substitution has not been
tested.

\paragraph{Torque floor.} We apply no floor to the torque, so above $\Rocrit$
the star's rotation freezes entirely. The Metcalfe et al.\ sample shows torques
20--70 times weaker than standard models rather than identically zero, so a
small residual torque may be more faithful. This affects the late-time rotation
periods in Figure~\ref{fig:rotation} but not the XUV, which is held up by its
own floor.

\paragraph{Exponent on the XUV.} We set $n = 1$ in
Equation~\ref{eq:apply}, applying the same suppression to the XUV as to the
torque. Since the torque depends on both the dipole field strength and the
mass-loss rate, and the X-ray luminosity enters only through the latter, a
value $n \ne 1$ is defensible; for dipole-dominated scalings anything between
$0.7$ and $1.4$ fits the Metcalfe et al.\ sample about equally well
($0.30$--$0.36$~dex scatter across that range). The choice matters at the
factor-of-a-few level for the late-time XUV and barely at all for cumulative
atmospheric escape, which is dominated by the saturated phase.

\begin{thebibliography}{99}

\bibitem[Baraffe et al.(2015)]{Baraffe15}
  Baraffe, I., Homeier, D., Allard, F., \& Chabrier, G. 2015, A\&A, 577, A42

\bibitem[Cameron \& Sch\"ussler(2017)]{Cameron17}
  Cameron, R.~H., \& Sch\"ussler, M. 2017, ApJ, 843, 111

\bibitem[Cranmer \& Saar(2011)]{CranmerSaar11}
  Cranmer, S.~R., \& Saar, S.~H. 2011, ApJ, 741, 54

\bibitem[Matt et al.(2015)]{Matt15}
  Matt, S.~P., Brun, A.~S., Baraffe, I., Bouvier, J., \& Chabrier, G. 2015,
  ApJL, 799, L23

\bibitem[Metcalfe et al.(2025)]{Metcalfe25}
  Metcalfe, T.~S., van Saders, J.~L., Pinsonneault, M.~H., et al. 2025, ApJL,
  accepted, arXiv:2509.03717

\bibitem[Ribas et al.(2005)]{Ribas05}
  Ribas, I., Guinan, E.~F., G\"udel, M., \& Audard, M. 2005, ApJ, 622, 680

\bibitem[van Saders et al.(2016)]{vanSaders16}
  van Saders, J.~L., Ceillier, T., Metcalfe, T.~S., et al. 2016, Nature, 529,
  181

\end{thebibliography}

\end{document}
